Place the red base at $z=0$ and its apex at $z=2$. At height $z$ its cross-section is the square $|x|,|y|\le a(z)$, where $a(z)=\sqrt2(1-\dfrac{z}{2})$. The blue cross-section is $|x|+|y|\le z$. On one red side, the visible horizontal length is $2a$ until $z_0=2\sqrt2-2$, then $4a-2z$ until $z_1=4-2\sqrt2$, then zero. A height increment multiplies this length by the face slant factor $\sqrt{\dfrac{3}{2}}$. Integrating the two linear widths gives $\int L(z)\,dz=z_0(\sqrt2+z_1)=10\sqrt2-12$. There are four red faces and four congruent blue contributions, hence lateral area $8\sqrt{\dfrac{3}{2}}(10\sqrt2-12)=80\sqrt3-48\sqrt6$. Both exposed bases have area eight. The total is $16+80\sqrt3-48\sqrt6$, so $a+b+c=\boxed{144}$.
\begin{asy}
import graph;import geometry;usepackage("amsmath");size(7cm);
real z=1,a=sqrt(2)*(1-z/2);draw(box((-a,-a),(a,a)),red);draw((z,0)--(0,z)--(-z,0)--(0,-z)--cycle,blue);label("red section",(0,-1.5),red);label("blue section at $z=1$",(0,1.5),blue);
\end{asy}

